\chapter{Agenci bazowi i infrastruktura eksperymentalna}
\label{chap:agenci-bazowi}

W niniejszym rozdziale przedstawiono projekt agentów bazowych oraz
infrastruktury służącej do uruchamiania i oceny strategii w środowisku
Ultimate Texas Hold'em. Strategie bazowe pełnią rolę punktów odniesienia
dla metod uczenia przez wzmacnianie. Pozwalają sprawdzić poprawność procesu
prowadzenia eksperymentów oraz określić, czy wyuczona polityka przewyższa
proste sposoby wyboru akcji.

Zaimplementowano wspólny kontrakt agenta, mechanizmy prowadzenia
pojedynczego epizodu i symulacji wielu rozdań, agenta losowego oraz
deterministycznego agenta regułowego. Infrastrukturę uzupełniono o
mechanizm agregacji metryk i zapis wyników eksperymentów. Wszystkie te
elementy korzystają z publicznego interfejsu środowiska opisanego w
rozdziale~\ref{chap:srodowisko-uth}.

\section{Wspólny interfejs agenta}
\label{sec:interfejs-agenta}

Wspólny kontrakt agentów zdefiniowano za pomocą protokołu
\texttt{Agent}. Każda zgodna implementacja udostępnia metodę
\texttt{select\_action()}, która przyjmuje obiekt
\texttt{UTHObservation} i zwraca jedną wartość typu \texttt{Action}.

% \begin{equation}
%     \operatorname{select\_action}
%     \colon
%     O_t
%     \rightarrow
%     a_t,
%     \qquad
%     a_t \in A_t^{legal}.
% \end{equation}

Zastosowanie protokołu pozwala korzystać z typowania strukturalnego.
Klasa agenta nie musi dziedziczyć po wspólnej klasie bazowej, lecz musi
udostępniać metodę o wymaganej sygnaturze. Dzięki temu ten sam mechanizm
prowadzenia rozgrywki może obsługiwać agentów losowych, regułowych
i uczących się bez znajomości ich wewnętrznej implementacji.

Agent otrzymuje wyłącznie obserwację zawierającą informacje dostępne
graczowi. Pełny obiekt \texttt{RoundState} nie jest częścią interfejsu,
dlatego implementacja strategii nie uzyskuje dostępu do kart krupiera,
kart spalonych ani kolejności pozostałych kart w talii.

\section{Prowadzenie pojedynczego epizodu}
\label{sec:runner-pojedynczego-epizodu}

Funkcja \texttt{play\_round()} łączy agenta z silnikiem
\texttt{UTHGame}. Rozpoczyna rozdanie, przekazuje bieżącą obserwację
agentowi, wykonuje wybraną akcję i powtarza ten proces do osiągnięcia
stanu terminalnego.

% Przepływ pojedynczego epizodu ma postać

% \begin{equation}
%     \operatorname{reset}
%     \rightarrow
%     O_0
%     \rightarrow
%     a_0
%     \rightarrow
%     \operatorname{step}
%     \rightarrow
%     \ldots
%     \rightarrow
%     S_T.
% \end{equation}

Każda decyzja zaakceptowana przez środowisko zostaje zapisana w kolejności
wykonania. Po zakończeniu rozdania funkcja zwraca obiekt
\texttt{EpisodeResult}, który zawiera sekwencję akcji, wynik rozdania
oraz pełne rozliczenie zakładów Ante, Blind i Play.

Model \texttt{EpisodeResult} udostępnia również wynik netto, łączną
wartość postawionych zakładów, liczbę decyzji i mnożnik zakładu Play.
Stanowi w ten sposób lekką reprezentację zakończonego epizodu, odpowiednią
do późniejszej agregacji wyników wielu rozdań.

\section{Agent losowy}
\label{sec:agent-losowy}

Agent losowy stanowi najprostszy punkt odniesienia. Nie analizuje wartości
kart ani prawdopodobieństwa wygranej. Dla każdej obserwacji wybiera jedną
z legalnych akcji.

% \begin{equation}
%     \pi_{\mathrm{random}}(a \mid O_t)
%     =
%     \frac{1}{\lvert A_t^{legal} \rvert},
%     \qquad
%     a \in A_t^{legal}.
% \end{equation}

Implementacja korzysta wyłącznie ze zbioru legalnych akcji zawartego
w \texttt{UTHObservation}. Akcje są porządkowane zgodnie z kolejnością
wartości enuma \texttt{Action}, dzięki czemu wynik generatora
pseudolosowego nie zależy od kolejności iterowania po zbiorze.

\subsection{Rozdzielenie źródeł losowości}
\label{subsec:rozdzielenie-losowosci}

Środowisko i agent posiadają niezależne generatory pseudolosowe. Generator
środowiska odpowiada za kolejność kart, natomiast generator agenta odpowiada
wyłącznie za wybór akcji:

\begin{equation}
    seed_{\mathrm{environment}}
    \rightarrow
    \text{kolejność kart},
    \qquad
    seed_{\mathrm{agent}}
    \rightarrow
    \text{wybór akcji}.
\end{equation}

Rozdzielenie obu źródeł losowości umożliwia zmianę strategii agenta bez
zmiany przebiegu rozdań. Pozwala również odtwarzać sekwencje decyzji
agenta niezależnie od konfiguracji środowiska.

\subsection{Walidacja implementacji}
\label{subsec:walidacja-agenta-losowego}

Testy agenta losowego obejmują zgodność ze wspólnym protokołem, wybieranie
wyłącznie legalnych akcji, powtarzalność decyzji dla tego samego seedu oraz
obsługę niepoprawnych danych wejściowych. Osobny test potwierdza, że agent
nie podejmuje decyzji dla obserwacji terminalnej.

Testy runnera obejmują wszystkie legalne ścieżki zakończenia rozdania.
Sprawdzają również, czy agent otrzymuje obserwacje odpowiadające kolejnym
fazom oraz czy jawnie przekazane źródło kart pozwala odtworzyć wynik
niezależnie od wewnętrznego seedu obiektu \texttt{UTHGame}.

\section{Agent regułowy}
\label{sec:agent-regulowy}

Drugą strategię bazową stanowi agent regułowy. W przeciwieństwie do
agenta losowego podejmuje on decyzje na podstawie kart widocznych
w obserwacji oraz aktualnej fazy rozdania. Przyjęte reguły oparto na
uproszczonej strategii Ultimate Texas Hold'em publikowanej przez
Wizard of Odds \cite{wizard_uth}.

Strategia została podzielona na trzy zbiory reguł odpowiadające
kolejnym punktom decyzyjnym środowiska.

\subsection{Decyzja przed flopem}

Przed odkryciem kart wspólnych agent wykonuje zakład Play równy
czterokrotności Ante dla następujących układów początkowych:

\begin{itemize}
    \item każdej pary od trójek wzwyż,
    \item każdego układu zawierającego asa,
    \item króla z dowolną kartą w tym samym kolorze,
    \item króla z kartą pięć lub wyższą w różnych kolorach,
    \item damy z kartą sześć lub wyższą w tym samym kolorze,
    \item damy z kartą osiem lub wyższą w różnych kolorach,
    \item waleta z kartą osiem lub wyższą w tym samym kolorze,
    \item waleta z dziesiątką w różnych kolorach.
\end{itemize}

Para dwójek oraz wszystkie pozostałe układy prowadzą do wykonania
akcji \texttt{CHECK}. Przyjęta strategia nie wykorzystuje zakładu
\texttt{BET\_3X}, mimo że akcja ta pozostaje legalną akcją środowiska.

\subsection{Decyzja po flopie}

Jeżeli gracz nie wykonał zakładu przed flopem, po odkryciu trzech
kart wspólnych agent wybiera \texttt{BET\_2X}, gdy posiada dwie pary
lub silniejszy układ, ukrytą parę z wyjątkiem pary dwójek albo cztery
karty do koloru z własną kartą tego koloru o randze co najmniej
dziesięć.

Przez ukrytą parę rozumiana jest para wykorzystująca co najmniej
jedną z dwóch prywatnych kart gracza. Jeżeli żaden z wymienionych
warunków nie jest spełniony, agent wykonuje akcję \texttt{CHECK}.

\subsection{Decyzja po odkryciu wszystkich kart wspólnych}

W ostatnim punkcie decyzyjnym agent wykonuje zakład
\texttt{BET\_1X}, jeżeli posiada ukrytą parę lub silniejszy układ.
Dla słabszych układów wykorzystywana jest reguła oparta na liczbie
kart mogących poprawić układ krupiera do układu pokonującego gracza.
Jeżeli liczba takich kart jest mniejsza niż 21, agent wykonuje
\texttt{BET\_1X}. W przeciwnym przypadku wybiera \texttt{FOLD}
\cite{wizard_uth}.

Tak zdefiniowana strategia jest deterministyczna. Dla tej samej
obserwacji agent zawsze zwraca tę samą akcję. Dzięki temu może
stanowić silniejszy punkt odniesienia niż agent losowy, jednocześnie
nie wykorzystując procesu uczenia ani informacji ukrytych przed
graczem.

\subsection{Implementacja strategii regułowej}
\label{subsec:implementacja-agenta-regulowego}

Reguły decyzyjne zostały rozdzielone od klasy realizującej wspólny
interfejs agenta. Funkcje \texttt{should\_raise\_preflop()},
\texttt{should\_raise\_flop()} oraz \texttt{should\_raise\_river()}
odpowiadają za ocenę obserwacji w poszczególnych fazach rozdania.
Klasa \texttt{RuleBasedAgent} wybiera na ich podstawie akcję zgodną
z aktualną fazą gry.

Dla decyzji na riverze zaimplementowano dodatkową funkcję
\texttt{count\_dealer\_outs()}. Tworzony jest zbiór kart niewidocznych
dla gracza, a następnie każda z nich jest analizowana jako pojedyncza
potencjalna karta prywatna krupiera. Jeżeli dołączenie danej karty do
kart wspólnych prowadzi do układu silniejszego od najlepszego układu
gracza, karta zostaje zaklasyfikowana jako out.

Podejście to odzwierciedla uproszczoną regułę 21 outów. Nie jest ono
pełnym przeszukaniem wszystkich możliwych dwukartowych układów
krupiera. Uwzględniane są pojedyncze niewidoczne karty, które same
prowadzą do układu pokonującego gracza. Pozwala to zachować charakter
strategii bazowej opisanej przez Wizard of Odds \cite{wizard_uth},
zamiast zastępować ją dokładniejszą metodą obliczeniową.

Klasa \texttt{RuleBasedAgent} nie przechowuje własnego stanu i nie
korzysta z generatora liczb pseudolosowych. Otrzymuje wyłącznie obiekt
\texttt{UTHObservation}, dzięki czemu nie posiada dostępu do kart
krupiera ani pozostałych informacji ukrytych w pełnym stanie
środowiska.

\section{Symulacja wielu epizodów}
\label{sec:symulacja-wielu-epizodow}

Do prowadzenia eksperymentów obejmujących wiele rozdań wykorzystano
obiekty \texttt{SimulationConfig} i \texttt{SimulationResult}.
Konfiguracja symulacji przechowuje uporządkowaną sekwencję seedów
wykorzystywanych do tworzenia talii. Liczba seedów określa jednocześnie
liczbę rozgrywanych epizodów.

Funkcja \texttt{run\_simulation()} dla każdego seedu tworzy osobną
talię oraz nową instancję środowiska \texttt{UTHGame}. Następnie
wykorzystuje funkcję \texttt{play\_round()} do rozegrania pełnego
epizodu. Wyniki wszystkich rozdań są przechowywane w obiekcie
\texttt{SimulationResult}.

Takie rozwiązanie oddziela konfigurację eksperymentu od implementacji
konkretnego agenta. Ten sam obiekt \texttt{SimulationConfig} może
zostać wykorzystany do oceny wielu różnych strategii.

\section{Powtarzalność eksperymentów}
\label{sec:powtarzalnosc-eksperymentow}

Kluczowym założeniem infrastruktury eksperymentalnej jest możliwość
odtworzenia przeprowadzonych symulacji. Dla każdego epizodu zapisywany
jest osobny seed talii. Ponowne utworzenie talii z tym samym seedem
prowadzi do tej samej kolejności kart.

Podczas porównywania agentów wykorzystywana jest identyczna sekwencja
seedów talii. Oznacza to, że każdy agent jest oceniany na tym samym
harmonogramie rozdań. Dzięki temu różnice wyników nie są skutkiem
wylosowania niezależnych zestawów kart dla poszczególnych strategii.

W przypadku agenta losowego dodatkowo wykorzystywany jest niezależny
seed odpowiedzialny wyłącznie za wybór akcji. Seed środowiska i seed
agenta nie wpływają więc na siebie. Agent regułowy jest deterministyczny
i nie wymaga własnego źródła losowości.

Konfiguracja wykorzystana w eksperymencie jest zapisywana razem
z wynikami. Obejmuje ona liczbę rozdań, seed generatora harmonogramu
talii, seed agenta losowego oraz pełną sekwencję seedów poszczególnych
talii. Pozwala to ponownie przeprowadzić dokładnie skonfigurowany
eksperyment.

\section{Metryki oceny}
\label{sec:metryki-agentow-bazowych}

Podstawową miarą jakości strategii jest wynik netto wyrażony
w jednostkach zakładu Ante. Dla symulacji składającej się z \(N\)
epizodów empiryczną wartość oczekiwaną oszacowano jako średni wynik
netto na jedno rozdanie:

\begin{equation}
    \widehat{EV}
    =
    \frac{1}{N}
    \sum_{i=1}^{N} r_i,
\end{equation}

gdzie \(r_i\) oznacza wynik netto \(i\)-tego epizodu wyrażony
w jednostkach Ante.

Oprócz średniego wyniku obliczane są całkowity wynik netto oraz
średnia i całkowita wartość postawionych zakładów. Zmienność wyników
opisano za pomocą odchylenia standardowego próby:

\begin{equation}
    s
    =
    \sqrt{
        \frac{
            \sum_{i=1}^{N}
            \left(r_i - \overline{r}\right)^2
        }{
            N - 1
        }
    }.
\end{equation}

Na jego podstawie wyznaczany jest błąd standardowy średniej:

\begin{equation}
    SE
    =
    \frac{s}{\sqrt{N}}.
\end{equation}

Infrastruktura zlicza również końcowe wyniki rozdań oraz częstotliwość
wykonywania poszczególnych akcji. Wyniki finansowe pojedynczych
epizodów i ich agregaty przechowywane są z wykorzystaniem typu
\texttt{Fraction}. Konwersja do liczb zmiennoprzecinkowych wykonywana
jest dopiero dla statystyk wymagających takiej reprezentacji oraz
podczas eksportu wyników.

\section{Porównanie agentów bazowych}
\label{sec:porownanie-agentow-bazowych}

W celu zweryfikowania kompletnego procesu eksperymentalnego
przeprowadzono porównanie agenta losowego i agenta regułowego.
Każdy agent rozegrał 10\,000 rozdań wykorzystujących ten sam
harmonogram talii. Seed generatora harmonogramu wynosił
\texttt{20260815}, natomiast generator decyzji agenta losowego
zainicjalizowano seedem \texttt{20260816}.

Uzyskane podstawowe metryki przedstawiono w
tabeli~\ref{tab:baseline-metrics}.

\begin{table}[ht]
    \centering
    \caption{Wyniki agentów bazowych dla 10\,000 rozdań}
    \label{tab:baseline-metrics}
    \footnotesize
    \begin{tabular}{lrrrrr}
        \hline
        Agent
        & $\widehat{EV}$
        & $s$
        & $SE$
        & Śr. stawka
        & Wynik netto \\
        \hline
        \texttt{RandomAgent}
        & $-0{,}46435$
        & $4{,}46422$
        & $0{,}04464$
        & $4{,}7395$
        & $-4643{,}5$ \\
        \texttt{RuleBasedAgent}
        & $0{,}03865$
        & $4{,}08130$
        & $0{,}04081$
        & $4{,}1905$
        & $386{,}5$ \\
        \hline
    \end{tabular}
\end{table}

Agent losowy uzyskał średni wynik
\(-0{,}46435\) jednostki Ante na rozdanie. Agent regułowy osiągnął
w tej samej próbie średni wynik \(0{,}03865\), czyli wynik wyższy
o około \(0{,}503\) jednostki Ante na rozdanie. Jednocześnie średnia
wartość środków stawianych przez agenta regułowego była niższa.

Dodatni wynik agenta regułowego w tej konkretnej próbie nie powinien
być interpretowany jako dowód dodatniej rzeczywistej wartości
oczekiwanej strategii. Błąd standardowy jego średniego wyniku wyniósł
około \(0{,}04081\), czyli wartość zbliżoną do samego oszacowania
\(\widehat{EV}\). Wynik eksperymentu służy przede wszystkim jako
punkt odniesienia dla późniejszej oceny agentów uczących się.

Rozkład wyników końcowych przedstawiono w
tabeli~\ref{tab:baseline-outcomes}.

\begin{table}[ht]
    \centering
    \caption{Rozkład wyników rozdań agentów bazowych}
    \label{tab:baseline-outcomes}
    \footnotesize
    \begin{tabular}{lrrrr}
        \hline
        Agent
        & Wygrana
        & Wygrana krupiera
        & Fold
        & Remis \\
        \hline
        \texttt{RandomAgent}
        & $45{,}01\%$
        & $42{,}75\%$
        & $8{,}49\%$
        & $3{,}75\%$ \\
        \texttt{RuleBasedAgent}
        & $47{,}92\%$
        & $31{,}75\%$
        & $16{,}79\%$
        & $3{,}54\%$ \\
        \hline
    \end{tabular}
\end{table}

Agent regułowy częściej kończył rozdania foldem, lecz znacznie rzadziej
doprowadzał do rozstrzygnięcia zakończonego wygraną krupiera.
Wyniki te pokazują, że zastosowane reguły prowadzą do jakościowo
innego sposobu podejmowania decyzji niż równomierny losowy wybór
spośród legalnych akcji.

\section{Zapis wyników eksperymentu}
\label{sec:zapis-wynikow-baseline}

Eksperyment porównujący agentów bazowych zapisuje wyniki w dwóch
formatach. Plik JSON zawiera konfigurację eksperymentu oraz
zagregowane metryki obu agentów. Plik CSV przechowuje dane
poszczególnych epizodów, w tym indeks rozdania, seed talii, nazwę
agenta, wynik netto, całkowitą stawkę, wynik końcowy i sekwencję
wykonanych akcji.

Zapis wyników na poziomie pojedynczego epizodu umożliwia późniejsze
analizowanie przebiegu eksperymentu bez konieczności ponownego
uruchamiania symulacji. Pozwala również zestawiać wyniki obu agentów
dla tych samych seedów talii oraz tworzyć dodatkowe statystyki
i wizualizacje.

\section{Walidacja infrastruktury eksperymentalnej}
\label{sec:walidacja-infrastruktury}

Testy infrastruktury obejmują zarówno pojedyncze epizody, jak
i symulacje wielu rozdań. Sprawdzana jest zgodność liczby epizodów
z konfiguracją, powtarzalność wyników dla tych samych seedów oraz
wykorzystanie tego samego harmonogramu talii przez porównywanych
agentów.

Testy agregacji metryk weryfikują obliczanie wyniku całkowitego,
średniego wyniku, wartości postawionych zakładów, odchylenia
standardowego, błędu standardowego oraz liczników wyników i akcji.
Osobne testy obejmują zapis wyników eksperymentu i sprawdzają,
czy zapisany harmonogram talii odpowiada konfiguracji symulacji oraz
czy dla obu agentów zapisano właściwą liczbę epizodów.

Takie rozdzielenie odpowiedzialności pozwala niezależnie testować
silnik gry, strategie agentów, prowadzenie epizodów, agregację metryk
oraz warstwę odpowiedzialną za zapis eksperymentów.

\section{Podsumowanie}
\label{sec:podsumowanie-agentow-bazowych}

W rozdziale zdefiniowano dwa punkty odniesienia dla dalszych
eksperymentów. \texttt{RandomAgent} reprezentuje strategię pozbawioną
wiedzy o wartości kart, natomiast \texttt{RuleBasedAgent} wykorzystuje
deterministyczny zestaw reguł oparty na uproszczonej strategii
Ultimate Texas Hold'em.

Przygotowana infrastruktura umożliwia prowadzenie powtarzalnych
symulacji na wspólnych harmonogramach rozdań, agregowanie wyników
i zapisywanie danych poszczególnych epizodów. Stanowi ona podstawę
do porównywania agentów wykorzystujących metody uczenia przez
wzmacnianie z ustalonymi strategiami bazowymi w kolejnych etapach
pracy.